\documentclass[12pt,a4paper]{article}

\usepackage[spanish,es-nodecimaldot]{babel}
\usepackage[utf8]{inputenc}
\usepackage[T1]{fontenc}
\usepackage{lmodern}
\usepackage{geometry}
\usepackage{microtype}
\usepackage{setspace}
\usepackage{parskip}
\usepackage{enumitem}
\usepackage{xcolor}
\usepackage{listings}
\usepackage{hyperref}
\usepackage{amsmath}
\usepackage{amssymb}
\usepackage{array}
\usepackage{booktabs}
\usepackage{longtable}

\geometry{margin=2.5cm}
\onehalfspacing
\setlist[itemize]{leftmargin=*,itemsep=0.2em}
\setlist[enumerate]{leftmargin=*,itemsep=0.2em}

\hypersetup{
  colorlinks=true,
  linkcolor=black,
  urlcolor=blue,
  pdftitle={Resolucion del Practico 3 - Testing de Software 2024},
  pdfauthor={}
}

\lstdefinestyle{codestyle}{
  basicstyle=\ttfamily\small,
  frame=single,
  breaklines=true,
  showstringspaces=false,
  keywordstyle=\color{blue!70!black}\bfseries,
  commentstyle=\color{gray!70!black}\itshape,
  stringstyle=\color{orange!70!black}
}

\title{\textbf{Resolucion del Practico 3}\\Testing de Software 2024}
\author{}
\date{}

\begin{document}

\maketitle
\tableofcontents
\newpage

\section*{Introduccion}
\addcontentsline{toc}{section}{Introduccion}

Este documento resuelve el \textit{Practico 3} de Testing de Software. El eje del practico es el
\textbf{Particionado del Espacio de Entradas} (\textit{Input Space Partitioning}, ISP) y los
criterios de cobertura combinatoria asociados (\textit{Each Choice}, \textit{Pair-Wise},
\textit{Base Choice}, \textit{All Combinations}). Toda la resolucion sigue el flujo recomendado por
Ammann y Offutt:

\begin{enumerate}
  \item identificar la unidad a testear y listar sus parametros,
  \item construir un \textit{Modelo del Dominio de Entradas} (MDE) a partir de la especificacion,
  \item declarar restricciones de factibilidad,
  \item aplicar un criterio para generar requisitos de test,
  \item instanciar valores concretos para satisfacer cada requisito.
\end{enumerate}

Cada ejercicio se acompana del MDE elegido, los requisitos de test (TR) derivados, los casos de
prueba que los cubren y, cuando corresponde, la implementacion que satisface dichos casos.

\section{Ejercicio 1: Lectura del Capitulo 6}

El ejercicio pide leer el Capitulo 6 (\textit{Input Space Partitioning}) del libro
\textit{Introduction to Software Testing} (Ammann y Offutt, 2da edicion). Los puntos centrales que
estructuran la lectura y que aparecen aplicados en los ejercicios siguientes son:

\begin{itemize}
  \item \textbf{Particionado del dominio.} Toda entrada valida debe caer en algun bloque
  (\textit{completitud}) y solo en uno (\textit{disjointness}). Si un valor cae en mas de un
  bloque o en ninguno, la particion esta mal definida y deben revisarse las caracteristicas.

  \item \textbf{Caracteristicas y bloques.} Una \textit{caracteristica} es un escenario relevante
  para testing (por ejemplo, ``la lista esta vacia''). Cada caracteristica induce una particion del
  dominio en \textit{bloques} (los posibles ``modos'' de esa propiedad).

  \item \textbf{Dos enfoques para construir el MDE.}
  \begin{itemize}
    \item \textbf{Basado en interfaz}: una caracteristica por parametro, mirando la sintaxis. Es
    rapido pero ignora relaciones entre parametros.
    \item \textbf{Basado en funcionalidad}: las caracteristicas se derivan del comportamiento
    esperado (precondiciones, poscondiciones, valores especiales). Es mas costoso pero mucho mas
    expresivo.
  \end{itemize}

  \item \textbf{Criterios de combinacion.}
  \begin{itemize}
    \item \textit{Each Choice Coverage (ECC)}: cada bloque aparece al menos en un test.
    \item \textit{Pair-Wise Coverage (PWC)}: todo par de bloques de caracteristicas distintas
    aparece en algun test.
    \item \textit{Base Choice Coverage (BCC)}: se fija un test base y se varia una caracteristica
    a la vez.
    \item \textit{All Combinations (ACoC)}: producto cartesiano de todos los bloques.
  \end{itemize}

  \item \textbf{Restricciones e infeasibility.} No todas las combinaciones son posibles. En ACoC y
  PWC se descartan las combinaciones invalidas; en BCC se reemplaza un valor por uno no base hasta
  obtener una combinacion factible.

  \item \textbf{Costo/beneficio.} PWC suele ser un buen compromiso entre tamano de la suite y
  capacidad de deteccion, ya que muchas fallas se disparan por interaccion de a dos factores.
\end{itemize}

Estos conceptos son la base directa de los ejercicios 2 a 5.

\section{Ejercicio 2: \texttt{numberOfOcurrences} -- PWC}

Se disena una bateria de tests para
\texttt{ListUtils.numberOfOcurrences(List<Integer> l, Integer element)} aplicando
\textit{Pair-Wise Coverage} sobre un MDE basado en funcionalidad.

\subsection{Archivos trabajados}

\begin{itemize}
  \item \texttt{src/main/java/assignment3\_exercises/ListUtils.java}
  \item \texttt{src/test/java/assignment3\_exercises/numberOfOcurrencesTest.java}
\end{itemize}

\subsection{Modelo del Dominio de Entradas (MDE)}

Se identifican cuatro caracteristicas relevantes a partir de la especificacion: dos verifican que
los parametros sean no nulos (precondicion explicita) y otras dos describen el comportamiento
funcional sobre la lista.

\begin{itemize}
  \item \textbf{C1 -- referencia de \texttt{l}}:
  \begin{itemize}
    \item L0: \texttt{l == null}
    \item L1: \texttt{l != null}
  \end{itemize}
  \item \textbf{C2 -- referencia de \texttt{element}}:
  \begin{itemize}
    \item E0: \texttt{element == null}
    \item E1: \texttt{element != null}
  \end{itemize}
  \item \textbf{C3 -- tamano de la lista} (solo aplica si L1):
  \begin{itemize}
    \item S0: lista vacia
    \item S1: lista no vacia
  \end{itemize}
  \item \textbf{C4 -- cantidad de ocurrencias} (solo aplica si L1 y E1):
  \begin{itemize}
    \item O0: 0 ocurrencias
    \item O1: 1 ocurrencia
    \item O2: mas de 1 ocurrencia
  \end{itemize}
\end{itemize}

\textbf{Restricciones de factibilidad}:
\begin{itemize}
  \item Si L0 entonces C3 y C4 no aplican: la lista nula dispara excepcion antes de inspeccionar
  cualquier otra propiedad.
  \item Si E0 entonces C4 no aplica: sin elemento valido no tiene sentido hablar de ocurrencias.
  \item Si S0 entonces necesariamente O0: una lista vacia no puede contener al elemento.
\end{itemize}

\subsection{Requisitos de test (PWC)}

PWC obliga a que toda pareja de bloques entre caracteristicas distintas aparezca en al menos un
test. Aplicando las restricciones anteriores, los pares factibles son:

\begin{itemize}[itemsep=0.05em]
  \item R01: (L0, E0) \quad R02: (L0, E1)
  \item R03: (L1, E0) \quad R04: (L1, E1)
  \item R05: (L1, S0) \quad R06: (L1, S1)
  \item R07: (L1, O0) \quad R08: (L1, O1) \quad R09: (L1, O2)
  \item R10: (E0, S0) \quad R11: (E0, S1)
  \item R12: (E1, S0) \quad R13: (E1, S1)
  \item R14: (E1, O0) \quad R15: (E1, O1) \quad R16: (E1, O2)
  \item R17: (S0, O0) \quad R18: (S1, O0) \quad R19: (S1, O1) \quad R20: (S1, O2)
\end{itemize}

\subsection{Casos de prueba}

\begin{longtable}{@{}p{1cm}p{4.5cm}p{2.5cm}p{2.2cm}p{3cm}@{}}
\toprule
\textbf{TC} & \textbf{Entrada \texttt{l}} & \textbf{\texttt{element}} & \textbf{Salida} & \textbf{Cubre} \\
\midrule
\endhead
TC1 & \texttt{null} & \texttt{null} & \texttt{IAE} & R01 \\
TC2 & \texttt{null} & 7 & \texttt{IAE} & R02 \\
TC3 & \texttt{[]} & \texttt{null} & \texttt{IAE} & R03, R05, R10 \\
TC4 & \texttt{[1,2,3]} & \texttt{null} & \texttt{IAE} & R03, R06, R11 \\
TC5 & \texttt{[]} & 5 & 0 & R04, R05, R07, R12, R14, R17 \\
TC6 & \texttt{[1,2,3]} & 9 & 0 & R04, R06, R07, R13, R14, R18 \\
TC7 & \texttt{[1,2,3]} & 2 & 1 & R04, R06, R08, R13, R15, R19 \\
TC8 & \texttt{[4,1,4,4]} & 4 & 3 & R04, R06, R09, R13, R16, R20 \\
\bottomrule
\end{longtable}

\noindent \textit{IAE} indica \texttt{IllegalArgumentException}. Adicionalmente se incluye un caso
de robustez (TC9) que verifica que la rutina no muta la lista de entrada: tras invocar
\texttt{numberOfOcurrences} sobre \texttt{[4,1,4,4]} buscando \texttt{4}, la lista debe seguir
siendo \texttt{[4,1,4,4]}.

\subsection{Implementacion}

\begin{lstlisting}[style=codestyle,language=Java]
public static int numberOfOcurrences(List<Integer> l, Integer element) {
    if (l == null || element == null) {
        throw new IllegalArgumentException("List and element must be non-null");
    }
    int count = 0;
    for (Integer current : l) {
        if (element.equals(current)) {
            count++;
        }
    }
    return count;
}
\end{lstlisting}

La rutina primero valida las precondiciones (R01--R03 dependen de esto) y luego recorre la lista
contando con \texttt{element.equals(current)} para que \texttt{Integer} se compare por valor y no
por referencia.

\subsection{Ejecucion}

\begin{lstlisting}[style=codestyle]
mvn -Dmaven.repo.local=.m2 -Dtest=numberOfOcurrencesTest test
\end{lstlisting}

\section{Ejercicio 3: \texttt{intersection} -- analisis del MDE y BCC}

El ejercicio analiza el MDE original propuesto en el enunciado para
\texttt{SetUtils.intersection(Set<Integer> set1, Set<Integer> set2)}, corrige sus problemas y
genera una suite por \textit{Base Choice Coverage}.

\subsection{Archivos trabajados}

\begin{itemize}
  \item \texttt{src/main/java/assignment3\_exercises/SetUtils.java}
  \item \texttt{src/test/java/assignment3\_exercises/IntersectionTest.java}
\end{itemize}

\subsection{a) Completitud de C1 (analogamente C2)}

\textbf{Si, C1 es completa.} Para un \texttt{Set<Integer> set1} solo hay tres posibilidades
mutuamente excluyentes:
\begin{itemize}
  \item \texttt{set1 == null},
  \item \texttt{set1 != null} y vacio,
  \item \texttt{set1 != null} con al menos un elemento.
\end{itemize}
Los bloques \texttt{b1}, \texttt{b2}, \texttt{b3} del enunciado cubren exactamente esas tres
posibilidades, por lo que todo \texttt{Set<Integer>} concebible cae en algun bloque.

\subsection{b) No solapamiento de C1 (analogamente C2)}

\textbf{Si, C1 no solapa.} Las tres situaciones son incompatibles entre si: un set nulo no puede
ser vacio o no vacio (ni siquiera tiene metodos invocables), y un set vacio no puede ser
simultaneamente no vacio. No existe \texttt{set1} que caiga a la vez en dos bloques.

\subsection{c) Completitud de C3}

\textbf{No, C3 no es completa.} Considerese el contraejemplo:
\[
\texttt{set1} = \{1,2\}, \qquad \texttt{set2} = \{2,3\}.
\]
Aqui hay interseccion no vacia (\(\{2\}\)), pero \texttt{set1} y \texttt{set2} no son iguales,
ninguno es subconjunto del otro y tampoco son disjuntos. Esta entrada no encaja en ninguno de los
bloques \texttt{b1}--\texttt{b4} originales: existe un \textit{solapamiento parcial} que el MDE no
contempla.

\subsection{d) Disjointness de C3}

\textbf{No, C3 no es disjunta.} Si ``subconjunto'' se interpreta de forma no estricta, el caso
\[
\texttt{set1} = \texttt{set2} = \{1,2\}
\]
cae simultaneamente en tres bloques: representan el mismo conjunto (\texttt{b1}), \texttt{set1} es
subconjunto de \texttt{set2} (\texttt{b2}) y \texttt{set2} es subconjunto de \texttt{set1}
(\texttt{b3}). Hay solapamiento entre los bloques.

\subsection{e) Revision del MDE}

Para que C3 sea completa y disjunta se utiliza la nocion de subconjunto \textit{propio} y se anade
un bloque adicional que cubra el solapamiento parcial:

\begin{itemize}
  \item \textbf{C1 (validez de \texttt{set1})}:
  \begin{itemize}
    \item C1.b1: \texttt{set1 == null}
    \item C1.b2: \texttt{set1 == \{\}}
    \item C1.b3: \texttt{set1} no vacio
  \end{itemize}
  \item \textbf{C2 (validez de \texttt{set2})}: analogo a C1.
  \item \textbf{C3 (relacion entre \texttt{set1} y \texttt{set2})}, solo aplicable si ninguno es
  nulo:
  \begin{itemize}
    \item C3.b1: \texttt{set1 = set2}
    \item C3.b2: \texttt{set1} es subconjunto \textit{propio} de \texttt{set2}
    \item C3.b3: \texttt{set2} es subconjunto \textit{propio} de \texttt{set1}
    \item C3.b4: \texttt{set1} y \texttt{set2} disjuntos
    \item C3.b5: solapamiento parcial (interseccion no vacia y ninguno es subconjunto del otro)
  \end{itemize}
\end{itemize}

Con estos ajustes C3 queda completa (todo par \((\texttt{set1},\texttt{set2})\) con ambos no nulos
encaja en alguno de los cinco bloques) y disjunta (los cinco bloques se excluyen mutuamente).

\subsection{f) Base Choice Coverage}

\textbf{Bloques base elegidos}: el caso ``camino feliz'' es ambos sets no vacios e iguales.
\begin{itemize}
  \item C1 base: C1.b3 (\texttt{set1} no vacio).
  \item C2 base: C2.b3 (\texttt{set2} no vacio).
  \item C3 base: C3.b1 (\texttt{set1 = set2}).
\end{itemize}

\textbf{Requisitos BC} (el test base mas un test por cada bloque no base, ajustando
factibilidad):

\begin{longtable}{@{}p{2cm}p{6cm}p{6.5cm}@{}}
\toprule
\textbf{TR} & \textbf{Variacion respecto del base} & \textbf{Resultado esperado} \\
\midrule
\endhead
TR-BC1 & --- (caso base) & \texttt{\{1,2,3\}} \\
TR-BC2 & C1 $\rightarrow$ C1.b1 (\texttt{null}) & \texttt{NullPointerException} \\
TR-BC3 & C1 $\rightarrow$ C1.b2 (\texttt{\{\}}) & \texttt{\{\}} \\
TR-BC4 & C2 $\rightarrow$ C2.b1 (\texttt{null}) & \texttt{NullPointerException} \\
TR-BC5 & C2 $\rightarrow$ C2.b2 (\texttt{\{\}}) & \texttt{\{\}} \\
TR-BC6 & C3 $\rightarrow$ C3.b2 (subc.\ propio) & \texttt{set1} \\
TR-BC7 & C3 $\rightarrow$ C3.b3 (superc.\ propio) & \texttt{set2} \\
TR-BC8 & C3 $\rightarrow$ C3.b4 (disjuntos) & \texttt{\{\}} \\
TR-BC9 & C3 $\rightarrow$ C3.b5 (solap.\ parcial) & elementos comunes \\
\bottomrule
\end{longtable}

\noindent Notar que TR-BC3 y TR-BC5 son los casos donde una variacion sobre C1 o C2 vuelve no
aplicable a C3; segun la guia de manejo de restricciones para BCC, se sustituye la base de C3 por
una combinacion factible (en este caso, interseccion vacia trivial).

\subsection{Casos implementados}

Cada \texttt{TC1}--\texttt{TC9} cubre uno de los TR-BCi correspondientes. Adicionalmente la suite
incluye:

\begin{itemize}
  \item Verificacion de no mutacion: \texttt{set1} y \texttt{set2} deben quedar igual a su copia
  pre-llamada.
  \item Verificacion de nueva instancia: el resultado no es \texttt{== set1} ni \texttt{== set2}.
  \item Borde adicional: ambos sets vacios $\rightarrow$ resultado vacio.
\end{itemize}

\subsection{Implementacion}

\begin{lstlisting}[style=codestyle,language=Java]
public static Set<Integer> intersection(Set<Integer> set1, Set<Integer> set2) {
    Objects.requireNonNull(set1, "set1 must be non-null");
    Objects.requireNonNull(set2, "set2 must be non-null");

    Set<Integer> result = new HashSet<>(set1);
    result.retainAll(set2);
    return result;
}
\end{lstlisting}

La validacion por \texttt{Objects.requireNonNull} respeta la especificacion de
\texttt{NullPointerException}, y la copia previa a \texttt{retainAll} garantiza tanto la no
mutacion de \texttt{set1} como la nueva instancia de retorno.

\subsection{Ejecucion}

\begin{lstlisting}[style=codestyle]
mvn -Dmaven.repo.local=.m2 -Dtest=IntersectionTest test
\end{lstlisting}

\section{Ejercicio 4: \texttt{PatternIndex} -- PWC y deteccion de fallas}

Se modela y prueba \texttt{PatternIndex.patternIndex(String subject, String pattern)} con
\textit{Pair-Wise Coverage}, y se documentan las fallas que la suite revela en la implementacion
original.

\subsection{Archivos trabajados}

\begin{itemize}
  \item \texttt{src/main/java/assignment3\_exercises/PatternIndex.java}
  \item \texttt{src/test/java/assignment3\_exercises/PatternIndexTest.java}
\end{itemize}

\subsection{Modelo del Dominio de Entradas}

El MDE combina dos caracteristicas de interfaz (referencia de cada parametro) con tres de
funcionalidad (tamano y posicion de la ocurrencia esperada):

\begin{itemize}
  \item \textbf{C1 -- referencia de \texttt{subject}}: C1.b1 = \texttt{null}, C1.b2 = no nulo.
  \item \textbf{C2 -- referencia de \texttt{pattern}}: C2.b1 = \texttt{null}, C2.b2 = no nulo.
  \item \textbf{C3 -- tamano de \texttt{subject}} (si C1.b2): C3.b1 = vacio, C3.b2 = no vacio.
  \item \textbf{C4 -- tamano de \texttt{pattern}} (si C2.b2): C4.b1 = vacio, C4.b2 = 1 caracter,
  C4.b3 = 2 o mas caracteres.
  \item \textbf{C5 -- posicion del resultado} (si C1.b2 y C2.b2): C5.b1 = indice 0 (inicio),
  C5.b2 = intermedio, C5.b3 = final, C5.b4 = \texttt{-1} (no encontrado).
\end{itemize}

\textbf{Restricciones}:
\begin{itemize}
  \item Si \texttt{subject} o \texttt{pattern} es \texttt{null}: precondicion violada,
  \texttt{IllegalArgumentException}; C3--C5 no aplican.
  \item Si \texttt{pattern} es vacio (C4.b1): por especificacion devuelve 0 (cubre C5.b1).
  \item Si \texttt{subject} es vacio (C3.b1) y \texttt{pattern} no lo es (C4.b2 o C4.b3): el
  resultado es \texttt{-1} (cubre C5.b4).
\end{itemize}

\subsection{Requisitos PWC y casos}

Los 20 requisitos de cobertura por pares (R01--R20) modelan combinaciones validas entre las
caracteristicas. La suite implementa 15 casos:

\begin{longtable}{@{}p{1cm}p{2.5cm}p{2.5cm}p{1.5cm}p{4cm}@{}}
\toprule
\textbf{TC} & \texttt{subject} & \texttt{pattern} & \textbf{Salida} & \textbf{Cubre} \\
\midrule
\endhead
TC1  & \texttt{null}     & \texttt{null}     & IAE & R01 \\
TC2  & \texttt{null}     & \texttt{"a"}      & IAE & R02 \\
TC3  & \texttt{"abc"}    & \texttt{null}     & IAE & R03 \\
TC4  & \texttt{""}       & \texttt{""}       & 0   & R04, R05, R11 \\
TC5  & \texttt{"abc"}    & \texttt{""}       & 0   & R04, R06, R12 \\
TC6  & \texttt{""}       & \texttt{"a"}      & -1  & R04, R07, R15 \\
TC7  & \texttt{"abc"}    & \texttt{"a"}      & 0   & R04, R08, R13 \\
TC8  & \texttt{"abc"}    & \texttt{"b"}      & 1   & R04, R08, R14 \\
TC9  & \texttt{"abc"}    & \texttt{"c"}      & 2   & R04, R08, R16 \\
TC10 & \texttt{"abcde"}  & \texttt{"ab"}     & 0   & R04, R09, R17 \\
TC11 & \texttt{"abcde"}  & \texttt{"cd"}     & 2   & R04, R09, R18 \\
TC12 & \texttt{"abcde"}  & \texttt{"de"}     & 3   & R04, R09, R19 \\
TC13 & \texttt{"abcde"}  & \texttt{"fg"}     & -1  & R04, R09, R20 \\
TC14 & \texttt{"abc"}    & \texttt{"z"}      & -1  & R04, R08, R15 \\
TC15 & \texttt{""}       & \texttt{"ab"}     & -1  & R04, R10, R20 \\
\bottomrule
\end{longtable}

\subsection{Fallas detectadas en la implementacion original}

La version provista en el repositorio \textit{no} sigue la especificacion documentada en el
Javadoc. Los tests revelan dos defectos concretos:

\begin{enumerate}
  \item \textbf{Ausencia de validacion de \texttt{null}.} La rutina invoca directamente
  \texttt{subject.length()} y \texttt{pattern.length()}, por lo que TC1--TC3 lanzan
  \texttt{NullPointerException} en lugar del \texttt{IllegalArgumentException} prometido.

  \item \textbf{Tratamiento incorrecto de \texttt{pattern} vacio.} Con
  \texttt{pattern.length() == 0}, la condicion \texttt{iSub + patternLen - 1 < subjectLen}
  evalua \(iSub - 1 < \text{subjectLen}\), lo cual es verdadero al menos para \(iSub = 0\). Dentro
  del cuerpo se ejecuta \texttt{pattern.charAt(0)}, que dispara
  \texttt{StringIndexOutOfBoundsException} en lugar de devolver 0.
\end{enumerate}

\subsection{Correccion aplicada}

Se anaden las dos guardas faltantes al principio del metodo:

\begin{lstlisting}[style=codestyle,language=Java]
public static int patternIndex(String subject, String pattern) {
    if (subject == null || pattern == null) {
        throw new IllegalArgumentException("subject and pattern must be non-null");
    }

    final int NOTFOUND = -1;
    int  iSub = 0, rtnIndex = NOTFOUND;
    boolean isPat  = false;
    int subjectLen = subject.length();
    int patternLen = pattern.length();

    if (patternLen == 0) {
        return 0;
    }

    while (isPat == false && iSub + patternLen - 1 < subjectLen) {
        if (subject.charAt(iSub) == pattern.charAt(0)) {
            rtnIndex = iSub;
            isPat = true;
            for (int iPat = 1; iPat < patternLen; iPat++) {
                if (subject.charAt(iSub + iPat) != pattern.charAt(iPat)) {
                    rtnIndex = NOTFOUND;
                    isPat = false;
                    break;
                }
            }
        }
        iSub++;
    }
    return rtnIndex;
}
\end{lstlisting}

Con estas dos correcciones, los 15 casos de la suite pasan: las precondiciones lanzan
\texttt{IllegalArgumentException}, el patron vacio devuelve 0 y los demas casos no se ven
afectados porque el algoritmo de busqueda original es correcto.

\subsection{Ejecucion}

\begin{lstlisting}[style=codestyle]
mvn -Dmaven.repo.local=.m2 -Dtest=PatternIndexTest test
\end{lstlisting}

\section{Ejercicio 5: \texttt{Iterator} sobre \texttt{ArrayList} -- PWC}

Se prueba la interfaz \texttt{java.util.Iterator} usando \texttt{ArrayList} como coleccion
subyacente (Iterator es interfaz y no se instancia directamente).

\subsection{Archivos trabajados}

\begin{itemize}
  \item \texttt{src/test/java/assignment3\_exercises/IteratorArrayListTest.java}
\end{itemize}

\subsection{Especificacion de referencia}

De Oracle Java SE 7 (\url{https://docs.oracle.com/javase/7/docs/api/java/util/Iterator.html}):

\begin{itemize}
  \item \texttt{hasNext()}: devuelve \texttt{true} si quedan elementos por iterar.
  \item \texttt{next()}: devuelve el proximo elemento; lanza \texttt{NoSuchElementException} si no
  hay mas.
  \item \texttt{remove()}: operacion opcional. Si no esta soportada, lanza
  \texttt{UnsupportedOperationException}. Si esta soportada, lanza \texttt{IllegalStateException}
  cuando no se llamo a \texttt{next()} previamente o cuando ya se llamo a \texttt{remove()}
  despues del ultimo \texttt{next()}.
\end{itemize}

\subsection{MDE}

\begin{itemize}
  \item \textbf{C1}: ``hay mas elementos'' $\in \{\text{true}, \text{false}\}$.
  \item \textbf{C2}: ``\texttt{next()} devuelve no nulo'' $\in \{\text{true}, \text{false}\}$.
  \item \textbf{C3}: ``\texttt{remove()} esta soportado'' $\in \{\text{true}, \text{false}\}$.
  \item \textbf{C4}: ``precondicion de \texttt{remove()} satisfecha''
  $\in \{\text{true}, \text{false}\}$.
\end{itemize}

\noindent \textbf{Decisiones de instanciacion}:
\begin{itemize}
  \item Para C3 = \texttt{true} se usa \texttt{new ArrayList<>(...).iterator()}, que sí soporta
  \texttt{remove()}.
  \item Para C3 = \texttt{false} se usa \texttt{Collections.unmodifiableList(...).iterator()},
  cuya implementacion delega en una vista que \textbf{no} soporta \texttt{remove()}, manteniendo
  un \texttt{ArrayList} como coleccion subyacente.
\end{itemize}

\subsection{Requisitos y casos}

Se cubren 13 requisitos:

\begin{longtable}{@{}p{1cm}p{2cm}p{9.5cm}@{}}
\toprule
\textbf{TC} & \textbf{R} & \textbf{Escenario} \\
\midrule
\endhead
TC1  & R01 & \texttt{hasNext()} = true sobre iterador con datos. \\
TC2  & R02 & \texttt{hasNext()} = false sobre iterador vacio. \\
TC3  & R03 & \texttt{hasNext()} = false sobre iterador agotado. \\
TC4  & R04 & \texttt{next()} devuelve un valor no nulo. \\
TC5  & R05 & \texttt{next()} devuelve \texttt{null} (la lista contiene un \texttt{null}). \\
TC6  & R06 & \texttt{next()} sobre iterador vacio $\rightarrow$ \texttt{NoSuchElementException}. \\
TC7  & R07 & \texttt{next()} sobre iterador agotado $\rightarrow$ \texttt{NoSuchElementException}. \\
TC8  & R08 & \texttt{remove()} no soportado (\texttt{unmodifiableList}) $\rightarrow$ \texttt{UOE}. \\
TC9  & R09 & \texttt{remove()} soportado pero sin \texttt{next()} previo $\rightarrow$ \texttt{ISE}. \\
TC10 & R10 & \texttt{remove()} sobre iterador vacio sin \texttt{next()} previo $\rightarrow$ \texttt{ISE}. \\
TC11 & R11 & \texttt{remove()} elimina el ultimo retornado con elementos restantes. \\
TC12 & R12 & \texttt{remove()} elimina el ultimo retornado sin elementos restantes. \\
TC13 & R13 & Doble \texttt{remove()} consecutivo $\rightarrow$ \texttt{ISE}. \\
\bottomrule
\end{longtable}

\subsection{Resultado}

La implementacion de \texttt{Iterator} para \texttt{ArrayList} (y la vista de
\texttt{unmodifiableList} para el caso C3 = false) respeta la especificacion en todos los
escenarios modelados. No se detectaron divergencias respecto del contrato documentado.

\subsection{Ejecucion}

\begin{lstlisting}[style=codestyle]
mvn -Dmaven.repo.local=.m2 -Dtest=IteratorArrayListTest test
\end{lstlisting}

\section*{Conclusiones}
\addcontentsline{toc}{section}{Conclusiones}

La resolucion consolida el trabajo con \textit{Input Space Partitioning} sobre cinco unidades
concretas. Tres lecciones se repiten a lo largo del practico:

\begin{enumerate}
  \item \textbf{El MDE no es opcional.} En el Ejercicio 3 quedo claro que un MDE descuidado (C3
  con bloques no disjuntos y no completos) hace inutil cualquier criterio que se aplique encima.
  Antes de generar requisitos hay que verificar completitud y disjointness.

  \item \textbf{PWC vs BCC.} PWC fuerza interacciones; BCC fuerza variaciones controladas
  alrededor de un camino feliz. La eleccion depende del problema: PWC es mas exhaustivo en
  interacciones (ejercicios 2 y 4), BCC concentra el esfuerzo en variaciones locales y se ajusta
  bien cuando se tiene un caso ``normal'' claro (ejercicio 3).

  \item \textbf{Trazabilidad.} Anclar cada test a un TR explicito hace que la suite sea auditable
  y revisable: se sabe \textit{por que} esta cada caso y que pasa si se elimina.
\end{enumerate}

\end{document}
